\chapter{The channel-ridge theorem: exact calibration by principal angles}

\section{Structured parental covariance and curvature}

The theorem is written in one fixed reference metric. Equivalently, first transform phenotype space by \(\M^{1/2}\) so that the chosen common metric becomes Euclidean, and then define ordinary orthogonal projectors in that transformed space. Let \(\Pp_i\) be the resulting rank-\(k\) projector onto population \(i\)'s channel. Define
\begin{align}
\G_i
&=g_{\perp}\Ii+(g_{\parallel}-g_{\perp})\Pp_i,
\label{eq:ridgeG}\\
\Hh_i
&=h_{\perp}\Ii-(h_{\perp}-h_{\parallel})\Pp_i.
\label{eq:ridgeH}
\end{align}
Set
\begin{equation}
\Delta g=g_{\parallel}-g_{\perp}>0,
\qquad
\Delta h=h_{\perp}-h_{\parallel}>0.
\end{equation}
Then
\begin{equation}
\G_i=g_{\perp}\Ii+\Delta g\Pp_i,
\qquad
\Hh_i=h_{\perp}\Ii-\Delta h\Pp_i.
\end{equation}

\begin{cautionbox}
The result is invariant to a change of trait units when the metric is transformed with those units, but it can depend on the scientifically chosen reference metric itself. In applications with different \(\G_i\), the projectors should therefore be constructed from the common-metric bases of \cref{eq:metricEuclideanization,eq:EuclideanizedChannelBasis}, not from raw population-specific whitened eigenvectors.
\end{cautionbox}

\inwords Genetic variance is enriched along the channel, while fitness curvature is reduced along the same channel. Each parental covariance is therefore coadapted to a shallow local ridge.

For \(\Pp_i\) to be the leading response subspace of \(\Sop_i\), require
\begin{equation}
g_{\parallel}h_{\parallel}>g_{\perp}h_{\perp}.
\label{eq:ridgeLeadingCondition}
\end{equation}
Indeed, along \(\Pp_i\) the response eigenvalue is \(g_{\parallel}h_{\parallel}\); in the complement it is \(g_{\perp}h_{\perp}\).

\section{Compute covariance and curvature differences}

Subtract the two covariance matrices:
\begin{equation}
\G_2-\G_1=\Delta g(\Pp_2-\Pp_1).
\end{equation}
Subtract the curvature matrices:
\begin{equation}
\Hh_2-\Hh_1=-\Delta h(\Pp_2-\Pp_1).
\end{equation}
With \(\Rr=\Zero\), substitute into \cref{eq:generalTraceIdentity}:
\begin{align}
L_{\mathrm{cov}}
&=-\frac18\tr\left[
-\Delta h(\Pp_2-\Pp_1)
\Delta g(\Pp_2-\Pp_1)
\right]\\
&=\frac{\Delta g\Delta h}{8}
\tr[(\Pp_2-\Pp_1)^2].
\end{align}
Use \cref{eq:chordalIdentity}:
\begin{equation}
\tr[(\Pp_2-\Pp_1)^2]
=\normF{\Pp_2-\Pp_1}^2
=2\sum_{j=1}^k\sin^2\psi_j.
\end{equation}
Therefore
\begin{theorem}[Channel-ridge fitness calibration]
Under \cref{eq:ridgeG,eq:ridgeH}, equal channel rank, zero segregation excess, and the parental benchmark above,
\begin{equation}
\boxed{
L_{\mathrm{cov}}
=\frac{\Delta g\Delta h}{4}
\sum_{j=1}^k\sin^2\psi_j
=\frac{\Delta g\Delta h}{4}d_C^2.}
\label{eq:ridgeTheorem}
\end{equation}
\end{theorem}

\inwords Channel mismatch becomes a positive fitness load because the exact same projector difference controls both covariance divergence and curvature divergence, with opposite signs.

\section{Early-time fitness snowball}

Combine \cref{eq:ridgeTheorem} with \cref{eq:localSnowballTheorem}:
\begin{align}
L_{\mathrm{cov}}(t)
&=\frac{\Delta g\Delta h}{4}
\left[t^2\normF{\cK}^2+O(t^3)\right]\\
&=\boxed{
\frac{\Delta g\Delta h}{4}\normF{\cK}^2t^2+O(t^3).}
\label{eq:fitnessSnowballLocal}
\end{align}

For equal principal angles \(\psi_j=\theta\),
\begin{equation}
L_{\mathrm{cov}}
=\frac{\Delta g\Delta h}{4}k\sin^2\theta
\approx\frac{\Delta g\Delta h}{4}k\theta^2.
\end{equation}

\begin{connectionbox}
The theory now has two distinct quadratic statements. Channel discrepancy is quadratic in time under deterministic forcing. Hybrid covariance load is proportional to channel discrepancy under covariance--curvature coadaptation. Only their combination justifies calling the result a reproductive-isolation snowball.
\end{connectionbox}

\begin{figure}[H]
\centering
\includegraphics[width=\textwidth]{figures/pedagogy_hybrid_load.pdf}
\caption{Geometric reading of the channel-ridge theorem. Parental genetic variance lies along differently oriented shallow ridges. The averaged hybrid covariance occupies directions that are steep in the parental fitness geometries.}
\label{fig:hybridLoadPedagogy}
\end{figure}

\begin{figure}[H]
\centering
\includegraphics[width=0.82\textwidth]{figures/03_hybrid_load_calibration.pdf}
\caption{Numerical calibration of raw channel geometry against the exact covariance--curvature load. Geometry acquires a fitness scale only after \(\Delta g\) and \(\Delta h\) are specified.}
\label{fig:hybridLoadCheck}
\end{figure}
